\clearpage\section*{Correction EV4-T2}\label{EV4-T2-CORR}
\textbf{QCM :} 1 D, 2 A, 3 B, 4 A, 5 B, 6 C, 7 A, 8 C, 9 C, 10 D, 11 D, 12 A. Chaque réponse correcte : 0,5 point ; total 6.
\paragraph{Vérifications des automatismes.} Priorités ; somme de fractions ; soustraction des exposants ; coefficient multiplicateur 1,2 ; division par 0,8 ; produit 1,25 fois 0,8 ; résolution affine ; accroissement vertical divisé par l’horizontal ; moyenne ; angle divisé par 360 ; complément à 1 ; substitution entre parenthèses.
\subsection*{Exercice 1 -- Un arbre renversé}
\begin{enumerate}
\item $P(A\cap B)=\dfrac{1}{5}$ ; $P(B)=\dfrac{1}{5}+\dfrac{3}{10}=\dfrac{1}{2}$. \textit{Barème : 1 point. Résultat et justification élémentaire attendus ; demi-point possible pour démarche correcte inachevée.}
\item $P_B(A)=\dfrac{2}{5}$ ; $P(A\cap\overline B)=\dfrac{2}{5}$ donc $P_{\overline B}(A)=\dfrac{4}{5}$. \textit{Barème : 2 points. Méthode ou justification : 1 ; calcul et conclusion : 1.}
\item Premier niveau $\dfrac{1}{2}$,$\dfrac{1}{2}$ ; après B : $\dfrac{2}{5}$ et $\dfrac{3}{5}$ ; après complément de B : $\dfrac{4}{5}$ et $\dfrac{1}{5}$. \textit{Barème : 2 points. Méthode ou justification : 1 ; calcul et conclusion : 1.}
\end{enumerate}
\subsection*{Exercice 2 -- Deux méthodes d’optimisation}
\begin{enumerate}
\item $f=16-(x-4)^2$ ; maximum 16 en 4. \textit{Barème : 1 point. Résultat et justification élémentaire attendus ; demi-point possible pour démarche correcte inachevée.}
\item $f\prime=8-2x$, positif avant 4 et négatif après. \textit{Barème : 2 points. Méthode ou justification : 1 ; calcul et conclusion : 1.}
\item Bornes : f(1)=7, f(6)=12 ; minimum 7. Inéquation $(x-4)^2\leq1$ : $[3;5]$. \textit{Barème : 2 points. Méthode ou justification : 1 ; calcul et conclusion : 1.}
\end{enumerate}
\subsection*{Exercice 3 -- Arguments et symétrie}
\begin{enumerate}
\item $A(-\dfrac{\sqrt 3}{2};-\dfrac{1}{2})$, $B(\dfrac{\sqrt 3}{2};-\dfrac{1}{2})$. \textit{Barème : 1 point. Résultat et justification élémentaire attendus ; demi-point possible pour démarche correcte inachevée.}
\item $-\dfrac{3}{4}+\dfrac{1}{4}=-\dfrac{1}{2}$. \textit{Barème : 1 point. Résultat et justification élémentaire attendus ; demi-point possible pour démarche correcte inachevée.}
\item Angle $\dfrac{2\pi}{3}$ ; $AB=\sqrt3$. \textit{Barème : 2 points. Méthode ou justification : 1 ; calcul et conclusion : 1.}
\end{enumerate}
\paragraph{Étalonnage.} Durée cible : 20–25 min pour les automatismes, 80–85 min pour les exercices, 10–15 min de relecture. Ces durées sont des estimations éditoriales à mesurer en classe. Accepter toute méthode correcte ; ne pas sanctionner deux fois une erreur numérique propagée avec cohérence. Le barème proposé ne vaut pas décision d’harmonisation de l’équipe.
